\documentclass[UTF8,zihao=-4]{ctexart}
\usepackage{amsmath,amssymb,geometry,enumitem,booktabs}
\geometry{a4paper,margin=2.2cm}
\setlist{nosep}
\title{第一章习题解答}
\author{按课堂板书方法整理}
\date{}
\begin{document}
\maketitle
\noindent\textbf{说明：}以下解答按课堂板书顺序书写：先写似然函数，再写先验分布，利用
\(\pi(\theta\mid x)\propto L(x\mid\theta)\pi(\theta)\) 得到后验核，最后补上归一化常数或识别后验分布。Gamma 分布采用课堂使用的率参数化。

\section*{1.1}
设不合格品数为 \(X\)，观测到 \(X=2\)，样本量 \(n=8\)。参数只可能取 \(0.1\) 和 \(0.2\)，且
\[
P(\theta=0.1)=0.7,\qquad P(\theta=0.2)=0.3.
\]

\textbf{Step 1：似然函数}
\[
P(X=2\mid\theta)=\binom{8}{2}\theta^2(1-\theta)^6.
\]
分别计算
\[
\begin{aligned}
P(X=2\mid\theta=0.1)&=\binom{8}{2}(0.1)^2(0.9)^6,\\
P(X=2\mid\theta=0.2)&=\binom{8}{2}(0.2)^2(0.8)^6.
\end{aligned}
\]

\textbf{Step 2：先验分布}
\[
P(\theta=0.1)=0.7,\qquad P(\theta=0.2)=0.3.
\]

\textbf{Step 3：后验分布}
由贝叶斯公式，
\[
P(\theta=0.1\mid X=2)
=\frac{P(X=2\mid\theta=0.1)P(\theta=0.1)}
{\sum_{j=1}^2P(X=2\mid\theta_j)P(\theta_j)}.
\]
组合数 \(\binom{8}{2}\) 在分子、分母中约去，故
\[
P(\theta=0.1\mid X=2)
=\frac{0.7(0.1)^2(0.9)^6}
{0.7(0.1)^2(0.9)^6+0.3(0.2)^2(0.8)^6}
\approx0.5418.
\]
同理，
\[
P(\theta=0.2\mid X=2)\approx0.4582.
\]
因此，\(\theta\) 的后验分布为
\[
\boxed{P(\theta=0.1\mid X=2)=0.5418,\qquad
P(\theta=0.2\mid X=2)=0.4582.}
\]

\section*{1.2}
设缺陷数 \(X\mid\lambda\sim\operatorname{Poisson}(\lambda)\)，参数只可能取 \(1.0,1.5\)，先验为
\[
P(\lambda=1.0)=0.4,\qquad P(\lambda=1.5)=0.6.
\]
观测到 \(X=3\)。

\textbf{Step 1：似然函数}
\[
P(X=3\mid\lambda)=\frac{e^{-\lambda}\lambda^3}{3!}.
\]

\textbf{Step 2：先验分布}
\[
P(\lambda=1.0)=0.4,\qquad P(\lambda=1.5)=0.6.
\]

\textbf{Step 3：后验分布}
\[
P(\lambda=1.0\mid X=3)
=\frac{0.4e^{-1}(1)^3/3!}
{0.4e^{-1}(1)^3/3!+0.6e^{-1.5}(1.5)^3/3!}.
\]
约去 \(3!\)，得
\[
P(\lambda=1.0\mid X=3)\approx0.2457.
\]
同理，
\[
P(\lambda=1.5\mid X=3)\approx0.7543.
\]
所以
\[
\boxed{P(\lambda=1.0\mid X=3)=0.2457,\qquad
P(\lambda=1.5\mid X=3)=0.7543.}
\]

\section*{1.3}
设从一批产品中抽取 \(n=8\) 件，发现不合格品数 \(x=3\)。
\[
X\mid\theta\sim\operatorname{Bin}(8,\theta),\qquad
L(\theta\mid x)=\binom{8}{3}\theta^3(1-\theta)^5.
\]

\subsection*{(1) \(\theta\sim U(0,1)\)}
\textbf{Step 1：似然函数}
\[
L(\theta\mid x)=\binom{8}{3}\theta^3(1-\theta)^5
\propto\theta^3(1-\theta)^5.
\]

\textbf{Step 2：先验密度}
\[
\pi(\theta)=1,\qquad0<\theta<1.
\]

\textbf{Step 3：后验核}
\[
\pi(\theta\mid x)\propto L(\theta\mid x)\pi(\theta)
\propto\theta^3(1-\theta)^5,\qquad0<\theta<1.
\]
这是 \(\operatorname{Be}(4,6)\) 的核。因此
\[
\boxed{\theta\mid x\sim\operatorname{Be}(4,6)},
\]
后验密度为
\[
\boxed{\pi(\theta\mid x)=\frac{\Gamma(10)}{\Gamma(4)\Gamma(6)}
\theta^3(1-\theta)^5=504\theta^3(1-\theta)^5,\quad0<\theta<1.}
\]

\subsection*{(2) \(\pi(\theta)=2(1-\theta),\ 0<\theta<1\)}
\textbf{Step 1：似然函数}
\[
L(\theta\mid x)\propto\theta^3(1-\theta)^5.
\]

\textbf{Step 2：先验密度}
\[
\pi(\theta)=2(1-\theta),\qquad0<\theta<1.
\]

\textbf{Step 3：后验核}
\[
\pi(\theta\mid x)\propto\theta^3(1-\theta)^5\cdot2(1-\theta)
\propto\theta^3(1-\theta)^6.
\]
这是 \(\operatorname{Be}(4,7)\) 的核。因此
\[
\boxed{\theta\mid x\sim\operatorname{Be}(4,7)},
\]
后验密度为
\[
\boxed{\pi(\theta\mid x)=\frac{\Gamma(11)}{\Gamma(4)\Gamma(7)}
\theta^3(1-\theta)^6=840\theta^3(1-\theta)^6,\quad0<\theta<1.}
\]

\section*{1.5}
设
\[
X\mid\theta\sim U\left(\theta-\frac12,\theta+\frac12\right),
\qquad\theta\sim U(10,20).
\]

\subsection*{(1) 观测到 \(X=12\)}
\textbf{Step 1：似然函数}
\[
p(12\mid\theta)=
\begin{cases}
1,&\theta-\frac12<12<\theta+\frac12,\\
0,&\text{其他情况}.
\end{cases}
\]
条件等价于 \(11.5<\theta<12.5\)，所以
\[
L(\theta\mid12)\propto I_{(11.5,12.5)}(\theta).
\]

\textbf{Step 2：先验密度}
\[
\pi(\theta)=\frac1{10},\qquad10<\theta<20.
\]

\textbf{Step 3：后验分布}
\[
\pi(\theta\mid12)\propto I_{(11.5,12.5)}(\theta)I_{(10,20)}(\theta)
=I_{(11.5,12.5)}(\theta).
\]
区间长度为 \(1\)，故
\[
\boxed{\theta\mid X=12\sim U(11.5,12.5)},
\qquad
\pi(\theta\mid12)=I_{(11.5,12.5)}(\theta).
\]

\subsection*{(2) 连续得到六个观测值}
六个观测值为
\[
12.0,\ 11.5,\ 11.7,\ 11.1,\ 11.4,\ 11.9.
\]
记 \(x=(x_1,\ldots,x_6)\)。由于观测独立，
\[
L(\theta\mid x)=\prod_{i=1}^6I_{(\theta-1/2,\theta+1/2)}(x_i).
\]
每个观测值都要求
\[
x_i-\frac12<\theta<x_i+\frac12.
\]
因此
\[
\max_i\left(x_i-\frac12\right)<\theta<\min_i\left(x_i+\frac12\right).
\]
本题中
\[
\max_i x_i=12.0,\qquad\min_i x_i=11.1,
\]
故
\[
11.5<\theta<11.6.
\]
与先验区间 \((10,20)\) 相交后仍为该区间，且其长度为 \(0.1\)。所以
\[
\boxed{\theta\mid x\sim U(11.5,11.6)},
\]
即
\[
\boxed{\pi(\theta\mid x)=10I_{(11.5,11.6)}(\theta).}
\]

\section*{1.6}
设 \(X_1,\ldots,X_n\mid\lambda\) 独立同分布于泊松分布 \(\operatorname{Poisson}(\lambda)\)，其中 \(\lambda\) 是均值参数。

\textbf{Step 1：似然函数}
\[
L(\lambda\mid x)=\prod_{i=1}^n\frac{e^{-\lambda}\lambda^{x_i}}{x_i!}
=\frac{e^{-n\lambda}\lambda^{\sum_i x_i}}{\prod_i x_i!}
\propto\lambda^{\sum_i x_i}e^{-n\lambda}.
\]

\textbf{Step 2：先验密度}
取 Gamma 先验 \(\lambda\sim\operatorname{Ga}(\alpha,\beta)\)，采用率参数化：
\[
\pi(\lambda)=\frac{\beta^\alpha}{\Gamma(\alpha)}
\lambda^{\alpha-1}e^{-\beta\lambda},\qquad\lambda>0.
\]

\textbf{Step 3：后验分布}
\[
\begin{aligned}
\pi(\lambda\mid x)&\propto L(\lambda\mid x)\pi(\lambda)\\
&\propto\lambda^{\alpha-1+\sum_i x_i}e^{-(\beta+n)\lambda}.
\end{aligned}
\]
这是 Gamma 分布的核，故
\[
\boxed{\lambda\mid x\sim\operatorname{Ga}\left(\alpha+\sum_{i=1}^nx_i,\ \beta+n\right).}
\]
因此 Gamma 分布是泊松均值参数 \(\lambda\) 的共轭先验分布。

\section*{1.7}
已知随机变量的密度为
\[
p(x\mid\theta)=\frac{2x}{\theta^2},\qquad0<x<\theta<1.
\]
设观测到 \(X=x\)，其中 \(0<x<1\)。

\subsection*{(1) \(\theta\sim U(0,1)\)}
\textbf{Step 1：似然函数}
\[
L(\theta\mid x)=\frac{2x}{\theta^2}I_{(x,1)}(\theta)
\propto\theta^{-2}I_{(x,1)}(\theta).
\]

\textbf{Step 2：先验密度}
\[
\pi(\theta)=1,\qquad0<\theta<1.
\]

\textbf{Step 3：后验分布}
\[
\pi(\theta\mid x)\propto\theta^{-2}I_{(x,1)}(\theta).
\]
归一化常数由
\[
\int_x^1\theta^{-2}\,d\theta=\frac1x-1=\frac{1-x}{x}
\]
给出，因此
\[
\boxed{\pi(\theta\mid x)=\frac{x}{1-x}\theta^{-2},\qquad x<\theta<1.}
\]

\subsection*{(2) \(\pi(\theta)=3\theta^2,\ 0<\theta<1\)}
\textbf{Step 1：似然函数}
\[
L(\theta\mid x)\propto\theta^{-2}I_{(x,1)}(\theta).
\]

\textbf{Step 2：先验密度}
\[
\pi(\theta)=3\theta^2,\qquad0<\theta<1.
\]

\textbf{Step 3：后验分布}
\[
\pi(\theta\mid x)\propto\theta^{-2}\cdot3\theta^2I_{(x,1)}(\theta)
\propto I_{(x,1)}(\theta).
\]
区间长度为 \(1-x\)，故
\[
\boxed{\pi(\theta\mid x)=\frac1{1-x},\qquad x<\theta<1,}
\]
即
\[
\boxed{\theta\mid x\sim U(x,1).}
\]

\section*{1.8}
抽取 \(n=100\) 件产品，发现不合格品数 \(x=3\)。先验为 \(\theta\sim\operatorname{Be}(2,200)\)。

\textbf{Step 1：似然函数}
\[
L(\theta\mid x)=\binom{100}{3}\theta^3(1-\theta)^{97}
\propto\theta^3(1-\theta)^{97}.
\]

\textbf{Step 2：先验密度}
\[
\pi(\theta)=\frac{\Gamma(202)}{\Gamma(2)\Gamma(200)}
\theta^{2-1}(1-\theta)^{200-1}
\propto\theta(1-\theta)^{199}.
\]

\textbf{Step 3：后验分布}
\[
\begin{aligned}
\pi(\theta\mid x)&\propto\theta^3(1-\theta)^{97}
\cdot\theta(1-\theta)^{199}\\
&\propto\theta^4(1-\theta)^{296}.
\end{aligned}
\]
所以
\[
\boxed{\theta\mid x\sim\operatorname{Be}(5,297).}
\]
其后验密度为
\[
\boxed{\pi(\theta\mid x)=\frac{\Gamma(302)}{\Gamma(5)\Gamma(297)}
\theta^4(1-\theta)^{296},\qquad0<\theta<1.}
\]

\section*{1.9}
设总体测量值
\[
X_i\mid\theta\overset{iid}{\sim}N(\theta,5^2),\qquad n=10,
\]
先验分布为
\[
\theta\sim N(172.72,2.54),
\]
这里按题目和课堂记号，第二个参数作为方差。样本均值为 \(\bar x=176.53\)。

\textbf{Step 1：似然函数}
因为 \(\bar X\mid\theta\sim N(\theta,5^2/10)\)，所以
\[
L(\theta\mid\bar x)\propto
\exp\left\{-\frac{n}{2\cdot5^2}(\bar x-\theta)^2\right\}.
\]

\textbf{Step 2：先验密度}
\[
\pi(\theta)\propto
\exp\left\{-\frac{(\theta-172.72)^2}{2\cdot2.54}\right\}.
\]

\textbf{Step 3：后验分布}
\[
\begin{aligned}
\pi(\theta\mid\bar x)&\propto
\exp\left\{-\frac12\left[\frac{10}{25}(\theta-176.53)^2
+\frac1{2.54}(\theta-172.72)^2\right]\right\}.
\end{aligned}
\]
这是正态分布的核。令
\[
A=\frac{10}{25}+\frac1{2.54}=0.7937008,
\]
则
\[
\tau_1^2=\frac1A=1.2599206,
\]
\[
\mu_1=\frac{(10/25)176.53+(1/2.54)172.72}{10/25+1/2.54}
=174.640119.
\]
所以
\[
\boxed{\theta\mid x\sim N(174.6401,1.2599).}
\]

\textbf{(2) 求区间}
正态后验密度关于后验均值 \(174.6401\) 对称，并且离均值越近密度越大。因此半径为 \(2.5\) 的所求区间为
\[
\boxed{[174.6401-2.5,\ 174.6401+2.5]
=[172.1401,177.1401].}
\]

\textbf{(3) 区间概率}
后验标准差为
\[
\tau_1=\sqrt{1.2599206}=1.1224619.
\]
因此
\[
\begin{aligned}
P(172.1401<\theta<177.1401\mid x)
&=P\left(-\frac{2.5}{1.1224619}<Z<\frac{2.5}{1.1224619}\right)\\
&=2\Phi(2.22725)-1\\
&\approx0.9741.
\end{aligned}
\]
\[
\boxed{P(\theta\text{ 落在该区间内}\mid x)\approx0.9741.}
\]

\section*{1.10}
设
\[
X_1,\ldots,X_n\mid\theta\overset{iid}{\sim}N(\theta,2^2),\qquad n=100,
\]
且先验 \(\theta\sim N(\mu,\tau^2)\)，其中 \(\tau^2>0\) 任意。

由正态--正态共轭结果，后验方差为
\[
\tau_1^2=\left(\frac n{\sigma^2}+\frac1{\tau^2}\right)^{-1}
=\left(\frac{100}{4}+\frac1{\tau^2}\right)^{-1}.
\]
因为 \(1/\tau^2>0\)，所以
\[
\tau_1^2<\left(\frac{100}{4}\right)^{-1}=\frac1{25}.
\]
开平方得后验标准差
\[
\boxed{\tau_1<\frac15.}
\]
这个结论与先验标准差 \(\tau\) 的具体大小无关，只要先验为正常的正态分布即可。

\section*{1.12}
设 \(X_1,\ldots,X_n\mid\theta\overset{iid}{\sim}U(0,\theta)\)，且先验为 Pareto 分布
\[
\pi(\theta)=
\begin{cases}
\displaystyle\frac{\alpha\theta_0^\alpha}{\theta^{\alpha+1}},&\theta>\theta_0,\\
0,&\theta\leq\theta_0,
\end{cases}
\qquad\theta_0>0,\ \alpha>0.
\]
记样本最大值为
\[
x_{(n)}=\max\{x_1,\ldots,x_n\}.
\]

\textbf{Step 1：似然函数}
单个观测的密度为
\[
p(x_i\mid\theta)=\frac1\theta I_{(0,\theta)}(x_i).
\]
因此样本似然为
\[
\begin{aligned}
L(\theta\mid x)&=\prod_{i=1}^n\frac1\theta I_{(0,\theta)}(x_i)\\
&=\theta^{-n}I(\theta\geq x_1,\ldots,\theta\geq x_n)\\
&=\theta^{-n}I(\theta\geq x_{(n)}).
\end{aligned}
\]
这里出现最大次序统计量，是因为 \(\theta\) 必须同时大于所有观测值，而这等价于只要求它大于最大观测值。

\textbf{Step 2：先验密度}
\[
\pi(\theta)=\frac{\alpha\theta_0^\alpha}{\theta^{\alpha+1}}
I(\theta>\theta_0).
\]

\textbf{Step 3：后验核}
\[
\begin{aligned}
\pi(\theta\mid x)&\propto L(\theta\mid x)\pi(\theta)\\
&\propto\theta^{-n}\theta^{-(\alpha+1)}
I(\theta\geq x_{(n)})I(\theta>\theta_0)\\
&\propto\theta^{-(\alpha+n+1)}I(\theta>a),
\end{aligned}
\]
其中
\[
a=\max\{\theta_0,x_{(n)}\}.
\]

\textbf{Step 4：归一化并识别分布}
因为
\[
\int_a^\infty\theta^{-(\alpha+n+1)}d\theta
=\frac{a^{-(\alpha+n)}}{\alpha+n},
\]
故后验密度为
\[
\boxed{\pi(\theta\mid x)=
(\alpha+n)a^{\alpha+n}\theta^{-(\alpha+n+1)}I(\theta>a).}
\]
这正是 Pareto 分布，因此
\[
\boxed{\theta\mid x\sim\operatorname{Pareto}(\alpha+n,a),
\qquad a=\max\{\theta_0,x_{(n)}\}.}
\]

\section*{1.15}
设 \(X_1,\ldots,X_n\mid\lambda\) 独立同分布于指数分布
\[
p(x\mid\lambda)=\lambda e^{-\lambda x},\qquad x>0.
\]

\subsection*{(1) 证明 Gamma 分布是共轭先验}
\textbf{Step 1：似然函数}
\[
L(\lambda\mid x)=\prod_{i=1}^n\lambda e^{-\lambda x_i}
=\lambda^n e^{-\lambda\sum_i x_i}.
\]

\textbf{Step 2：先验密度}
取 \(\lambda\sim\operatorname{Ga}(\alpha,\beta)\)，即
\[
\pi(\lambda)=\frac{\beta^\alpha}{\Gamma(\alpha)}
\lambda^{\alpha-1}e^{-\beta\lambda},\qquad\lambda>0.
\]

\textbf{Step 3：后验分布}
\[
\begin{aligned}
\pi(\lambda\mid x)&\propto L(\lambda\mid x)\pi(\lambda)\\
&\propto\lambda^n e^{-\lambda\sum_i x_i}
\lambda^{\alpha-1}e^{-\beta\lambda}\\
&\propto\lambda^{\alpha+n-1}
e^{-\left(\beta+\sum_i x_i\right)\lambda}.
\end{aligned}
\]
所以
\[
\boxed{\lambda\mid x\sim\operatorname{Ga}\left(\alpha+n,\ \beta+\sum_{i=1}^nx_i\right).}
\]
后验仍为 Gamma 分布，故 Gamma 分布是参数 \(\lambda\) 的共轭先验分布。

\subsection*{(2) 确定超参数}
Gamma 分布的均值和方差为
\[
E(\lambda)=\frac\alpha\beta,\qquad
\operatorname{Var}(\lambda)=\frac\alpha{\beta^2}.
\]
题给先验均值为 \(0.0002\)，先验标准差为 \(0.0001\)，所以方差为 \(10^{-8}\)。于是
\[
\frac\alpha\beta=0.0002,\qquad
\frac\alpha{\beta^2}=10^{-8}.
\]
两式相除得
\[
\beta=\frac{0.0002}{10^{-8}}=20000,
\]
再由第一式得
\[
\alpha=0.0002\times20000=4.
\]
故先验分布为
\[
\boxed{\lambda\sim\operatorname{Ga}(4,20000).}
\]

\section*{1.19}
设 \(X_1,\ldots,X_n\mid\lambda\) 是来自泊松分布 \(\operatorname{Poisson}(\lambda)\) 的样本，证明
\[
T(X)=\sum_{i=1}^nX_i
\]
是 \(\lambda\) 的充分统计量。

样本的联合密度（即似然函数）为
\[
\begin{aligned}
L(\lambda\mid x)&=\prod_{i=1}^n\frac{e^{-\lambda}\lambda^{x_i}}{x_i!}\\
&=e^{-n\lambda}\lambda^{\sum_{i=1}^nx_i}
\frac1{\prod_{i=1}^nx_i!}.
\end{aligned}
\]
令
\[
T(x)=\sum_{i=1}^nx_i,qquad
g(T(x),\lambda)=e^{-n\lambda}\lambda^{T(x)},qquad
h(x)=\frac1{\prod_{i=1}^nx_i!}.
\]
则
\[
L(\lambda\mid x)=g(T(x),\lambda)h(x).
\]
其中 \(g\) 只通过统计量 \(T(x)\) 和参数 \(\lambda\) 发生变化，而 \(h(x)\) 与 \(\lambda\) 无关。根据因子分解定理，
\[
\boxed{T(X)=\sum_{i=1}^nX_i\text{ 是 }\lambda\text{ 的充分统计量}.}
\]
\end{document}
